\documentclass[../../../include/open-logic-section]{subfiles}

\begin{document}

\olfileid{sth}{choice}{vitali}
\olsection{پیوست: تناقض‌نمای ویتالی}

برای آن‌که واقعاً دریابیم ساخت باناخ--تارسکی پذیرفتنی است یا نه،
باید \emph{اثبات} آن را بررسی کنیم. متأسفانه این کار به جبر بسیار
بیشتری از آنچه می‌توانیم اینجا عرضه کنیم نیاز دارد. بااین‌حال،
می‌توانیم چند نکتهٔ کوتاه مطرح کنیم که شاید با تمرکز بر نتیجهٔ زیر
پرتوی بر اثبات باناخ--تارسکی بیفکند:\footnote{برای شرحی بسیار
کامل‌تر، نگاه کنید به \cite{Weston2003} یا \cite{Wagon2016}.}

\begin{thm}[تناقض‌نمای ویتالی (در $\ZFC$)]\ollabel{vitaliparadox}
هر دایره را می‌توان به شمارایی قطعه تجزیه کرد و آن قطعه‌ها را (با
دوران و انتقال) دوباره چنان کنار هم گذاشت که دو نسخه از همان دایره
ساخته شود.
\end{thm}

اثبات تناقض‌نمای ویتالی بسیار آسان‌تر از تناقض‌نمای باناخ--تارسکی
است. آن را «تناقض‌نمای ویتالی» نامیده‌ایم، زیرا از ساخت یک مجموعهٔ
اندازه‌ناپذیر به دست ویتالی در \citeyear{Vitali1905} نتیجه می‌شود.
اما جنبه‌های نظریه‌مجموعه‌ایِ اثبات تناقض‌نمای ویتالی و تناقض‌نمای
باناخ--تارسکی بسیار شبیه‌اند. تفاوت اساسیِ نتایج فقط این است که
باناخ--تارسکی تجزیه‌ای \emph{متناهی} را در نظر می‌گیرد، حال آن‌که
تناقض‌نمای ویتالی تجزیه‌ای \emph{نامتناهیِ شمارا} را بررسی می‌کند.
به تعبیر \citet{Weston2003}، تناقض‌نمای ویتالی «بی‌تردید به‌اندازهٔ
تناقض‌نمای باناخ--تارسکی چشمگیر نیست، اما نشان می‌دهد که تناقض‌نماهای
هندسی حتی در وضعیت‌های ساده نیز می‌توانند رخ دهند».

تناقض‌نمای ویتالی دربارهٔ شکلی دوبعدی، یعنی یک دایره، است. پس روی
صفحهٔ $\Real^2$ کار می‌کنیم. بگذارید $\rotationsgroup$ مجموعهٔ
دوران‌های (ساعت‌گردِ) نقاط حول مبدأ با زاویه‌های
$2\pi q$ باشد، که در آن
$q \in \mathbb{Q}\cap[0,1)$. به بیان دیگر، زاویه‌ها مضرب‌های گویا
از یک دور کامل‌اند. چند واقعیت جبری دربارهٔ $\rotationsgroup$ در
زیر آمده است (اگر صورت نتیجه را نمی‌فهمید، اثبات معنای آن را روشن
خواهد کرد):

\begin{lem}\ollabel{rotationsgroupabelian}
$\rotationsgroup$ تحت ترکیب توابع یک {گروهِ} آبلی تشکیل می‌دهد.
\end{lem}

\begin{proof}
$0_{\rotationsgroup}$ را برای دوران به اندازهٔ $0$ رادیان بنویسید.
این عضو همانیِ $\rotationsgroup$ است، زیرا برای هر
$\rho \in \rotationsgroup$:
$\comp{0_{\rotationsgroup}}{\rho} =
\comp{\rho}{0_{\rotationsgroup}} = \rho$.

هر عضو وارون دارد. اگر $\rho \in \rotationsgroup$ به اندازهٔ $r$
رادیان دوران دهد، $\rho^{-1} \in \rotationsgroup$ به اندازهٔ
$2\pi-r$ رادیان دوران می‌دهد؛ پس
$\rho \circ \rho^{-1} = 0_{\rotationsgroup}$.

ترکیب شرکت‌پذیر است: برای هر
$\rho, \sigma, \tau \in \rotationsgroup$،
$\comp{\rho}{(\comp{\sigma}{\tau})} =
\comp{(\comp{\rho}{\sigma})}{\tau}$.

ترکیب جابجایی‌پذیر است: برای هر
$\rho, \sigma \in \rotationsgroup$،
$\comp{\rho}{\sigma} = \comp{\sigma}{\rho}$.
\end{proof}

در واقع، می‌توانیم گروه $\rotationsgroup$ را به دو نیم تقسیم کنیم و
سپس با هریک از آن دو نیم، کل گروه را بازیابیم:

\begin{lem}\ollabel{disjointgroup}
افرازی از $\rotationsgroup$ به دو مجموعهٔ مجزای
$\rotationsgroup_{1}$ و $\rotationsgroup_{2}$ وجود دارد، به‌طوری‌که
برای $i=1,2$:
\[
\Setabs{\comp{\rho}{\rho}}{\rho \in \rotationsgroup_i}
= \rotationsgroup.
\]
\end{lem}

\begin{proof}
بگذارید $\rotationsgroup_{1}$ از دوران‌های $\rotationsgroup$ با
زاویه‌هایی در $[0,\pi)$ تشکیل شود؛ و بگذارید
$\rotationsgroup_{2} = \rotationsgroup
\setminus \rotationsgroup_1$. بنا بر جبر مقدماتی،
$\Setabs{\comp{\rho}{\rho}}{\rho \in \rotationsgroup_1} =
\rotationsgroup$. نتیجه‌ای مشابه برای $\rotationsgroup_2$ به دست
می‌آید.
\end{proof}

از این واقعیت دربارهٔ گروه‌ها برای اثبات
\olref{vitaliparadox} استفاده می‌کنیم. بگذارید $\onesphere$ دایرهٔ
واحد باشد؛ یعنی مجموعهٔ نقاطی که دقیقاً~$1$ واحد از مبدأ صفحه فاصله
دارند، یعنی
$\Setabs{\tuple{r,s} \in \Real^2}{\sqrt{r^2+s^2}=1}$. با در نظر
گرفتن رابطهٔ زیر روی $\onesphere$، آن را به بخش‌هایی تقسیم می‌کنیم:
\[
	r \sim s \emph{ اگر و تنها اگر }(\exists \rho \in \rotationsgroup)\rho(r) = s.
\]
یعنی نقاط $\onesphere$ تحت این رابطه با هم مرتبط‌اند اگر و تنها اگر
بتوان با دورانی حول مبدأ که زاویه‌اش مضربی گویا از یک دور کامل است،
از یکی به دیگری رسید. چنان‌که انتظار می‌رود:

\begin{lem}
$\sim$ یک رابطهٔ هم‌ارزی است.
\end{lem}

\begin{proof}
با استفاده از \olref{rotationsgroupabelian} بدیهی است.
\end{proof}

اکنون اصل انتخاب را به کار می‌بریم تا مجموعه‌ای مانند $C$ به دست
آوریم که دقیقاً یک عضو از هر کلاس هم‌ارزیِ $\onesphere$ تحت
$\sim$ دارد. یعنی تابع انتخابی مانند $f$ را روی مجموعهٔ کلاس‌های
هم‌ارزی در نظر می‌گیریم،\footnote{چون $\rotationsgroup$
!!{enumerable} است، هر !!{element} از $E$ نیز !!{enumerable} است.
چون $\onesphere$ !!{nonenumerable} است، از
\olref{vitalicover} و
\olref[card-arithmetic][simp]{kappaunionkappasize} نتیجه می‌شود که
$E$ !!{nonenumerable} است. پس این کاربردی از انتخاب
\emph{نا}شماراست.}
\[
	E = \Setabs{\equivrep{r}{\sim}}{r \in \onesphere},
\]
و می‌گذاریم $C = \ran{f}$. برای هر دوران
$\rho \in \rotationsgroup$، مجموعهٔ $\funimage{\rho}{C}$ از نقاطی
تشکیل می‌شود که با اعمال دوران $\rho$ به هر نقطهٔ $C$ به دست
می‌آیند. دو نتیجهٔ بعدی نشان می‌دهند که این مجموعه‌ها دایره را
به‌طور کامل و بدون هم‌پوشانی می‌پوشانند:

\begin{lem}\ollabel{vitalicover}
$\onesphere = \bigcup_{\rho \in \rotationsgroup} \funimage{\rho}{C}$.
\end{lem}

\begin{proof}
$s \in \onesphere$ را ثابت کنید؛ یک $r \in C$ هست به‌طوری‌که
$r \in \equivrep{s}{\sim}$؛ یعنی $r \sim s$، و بنابراین برای یک
$\rho \in \rotationsgroup$ داریم $\rho(r) = s$.
\end{proof}

\begin{lem}\ollabel{vitalinooverlap}
اگر $\rho_1 \neq \rho_2$، آن‌گاه
$\funimage{\rho_1}{C} \cap \funimage{\rho_2}{C} = \emptyset$.
\end{lem}

\begin{proof}
فرض کنید
$s \in \funimage{\rho_{1}}{C} \cap \funimage{\rho_{2}}{C}$. پس برای
یک $r_{1}, r_{2} \in C$ داریم
$s = \rho_{1}(r_{1}) = \rho_{2}(r_{2})$. بنابراین
$\rho^{-1}_2(\rho_1(r_1)) = r_2$ و
$\comp{\rho_1}{\rho^{-1}_2} \in \rotationsgroup$؛ پس
$r_{1} \sim r_{2}$. چون $C$ دقیقاً یک عضو از هر کلاس هم‌ارزی تحت
$\sim$ برمی‌گزیند، $r_1 = r_2$. پس
$s = \rho_1(r_1) = \rho_2(r_1)$ و در نتیجه
$\rho_1 = \rho_2$.
\end{proof}

اکنون واقعیت‌های جبریِ پیشین را بر دایره اعمال می‌کنیم:

\begin{lem}\ollabel{pseudobanachtarski}
افرازی از $\onesphere$ به دو مجموعهٔ مجزای $D_{1}$ و $D_{2}$ وجود
دارد، به‌طوری‌که می‌توان $D_{1}$ را به شمارایی مجموعه افراز کرد و
آن‌ها را با دوران به نسخه‌ای از $\onesphere$ تبدیل کرد (و برای
$D_{2}$ نیز همین‌طور).
\end{lem}

\begin{proof}
با استفاده از $\rotationsgroup_{1}$ و $\rotationsgroup_{2}$ در
\olref{disjointgroup}، بگذارید:
\begin{align*}
	D_{1} &= \bigcup_{\rho \in \rotationsgroup_1} \funimage{\rho}{C} &
	D_{2} &= \bigcup_{\rho \in \rotationsgroup_2} \funimage{\rho}{C}
\end{align*}
بنا بر \olref{vitalicover}، این افرازی از $\onesphere$ است؛ و بنا
بر \olref{vitalinooverlap}، $D_1$ و $D_2$ مجزایند. بنا بر ساخت،
$D_1$ را می‌توان به شمارایی مجموعه، یعنی
$\funimage{\rho}{C}$ برای هر
$\rho \in \rotationsgroup_1$، افراز کرد. این مجموعه‌ها را می‌توان
با دوران به نسخه‌ای از $\onesphere$ تبدیل کرد، زیرا بنا بر
\olref{disjointgroup} و \olref{vitalicover}:
$\onesphere = \bigcup_{\rho \in
\rotationsgroup}\funimage{\rho}{C} = \bigcup_{\rho \in
\rotationsgroup_1}\funimage{(\comp{\rho}{\rho})}{C}$.
همین استدلال دربارهٔ $D_2$ برقرار است.
\end{proof}\noindent
این بی‌درنگ تناقض‌نمای ویتالی را نتیجه می‌دهد. زیرا فقط با تقسیم
$\onesphere$ به شمارایی قطعه (مجموعه‌های گوناگون
$\funimage{\rho}{C}$) و سپس حرکت صلب آن‌ها، می‌توانیم از
$\onesphere$ \emph{دو} نسخه بسازیم: هر قطعهٔ $D_1$ را دوران دهید،
و هر قطعهٔ $D_2$ را نخست انتقال و سپس دوران دهید.

راهبرد اثبات را مرور کنیم. با چند واقعیت جبری دربارهٔ گروه دوران‌ها
روی صفحه آغاز کردیم. از این گروه برای افراز $\onesphere$ به
کلاس‌های هم‌ارزی استفاده کردیم. سپس با استفاده از انتخاب برای
برگزیدن اعضایی از هر کلاس، به یک «تناقض‌نما» رسیدیم.

برای اثبات باناخ--تارسکی دقیقاً از همین راهبرد استفاده می‌کنیم.
تفاوت اصلی این است که واقعیت‌های جبریِ مورد استفاده در اثبات
باناخ--تارسکی به‌مراتب پیچیده‌تر از واقعیت‌های لازم برای
تناقض‌نمای ویتالی‌اند. اما آن واقعیت‌های جبری ربطی به انتخاب ندارند.
آن‌ها را به‌اختصار خلاصه می‌کنیم.

برای اثبات باناخ--تارسکی، با همتایی برای
\olref{disjointgroup} آغاز می‌کنیم: هر \emph{گروه آزاد} را می‌توان
به چهار قطعه تقسیم کرد که از نظر شهودی می‌توان آن‌ها را «جابجا کرد»
تا دو نسخه از کل گروه بازیابی شود.\footnote{امکان استفاده از
\emph{چهار} قطعه نتیجهٔ \cite{Robinson1947} است. برای اثباتی جدید،
نگاه کنید به \citet[Theorem 5.2]{Wagon2016}. در توصیف این فرایند
به‌صورت «جابجاکردن» قطعات گروه از
\citet[p.~3]{Weston2003} پیروی می‌کنیم.} سپس نشان می‌دهیم که با دو
دوران مشخص حول مبدأ $\Real^3$ می‌توان گروه آزادی از دوران‌ها، یعنی
$F$، تولید کرد.\footnote{نگاه کنید به
\citet[Theorem 2.1]{Wagon2016}.} (هنوز انتخاب به کار نرفته است.)

عملِ $F$ روی سطح کره آزاد نیست: هر دوران نابدیهی نقاط ثابتی دارد.
پس نخست مجموعهٔ شمارای نقاطی را کنار می‌گذاریم که عضوی نابدیهی از
$F$ ثابتشان می‌کند. روی باقیِ سطح، نقاط را «مشابه» می‌گیریم اگر و
تنها اگر یکی با دورانی از~$F$ از دیگری به دست آید. سپس
\emph{انتخاب را به کار می‌بریم} تا دقیقاً یک نقطه از هر کلاس
هم‌ارزیِ نقاط «مشابه» برگزینیم. با اعمال تقسیم $F$ به این بخشِ سطح،
همانند \olref{pseudobanachtarski}، آن را به چهار قطعه می‌شکنیم که
می‌توان «جابجایشان کرد» و دو نسخه از آن بخش ساخت. سرانجام، با گام
استانداردِ جذب، مجموعهٔ شمارای کنارگذاشته‌شده را دوباره به قطعات
می‌افزاییم و کل سطح کره را بازیابی می‌کنیم. این کار
\citep{Hausdorff1914} را اثبات می‌کند:

\begin{thm}[تناقض‌نمای هاوسدورف (در $\ZFC$)]
سطح هر کره را می‌توان به شمار متناهی‌ای قطعه تجزیه کرد و آن قطعه‌ها
را (با دوران و انتقال) دوباره چنان کنار هم گذاشت که دو نسخهٔ مجزا
از همان کره ساخته شود.
\end{thm}

برای به‌دست‌آوردن قضیهٔ کامل باناخ--تارسکی (که نه‌فقط سطح کره، بلکه
درون آن را نیز دربر می‌گیرد) به چند ترفند جبری دیگر نیاز است. اما
راستش را بخواهید، این دیگر فقط روکشِ کیک جبری است. ازاین‌رو وستون
می‌نویسد:
\begin{quote}
  [\ldots] نتیجه دربارهٔ گروه‌های آزاد \emph{گام کلیدی} در اثبات
  تناقض‌نمای باناخ--تارسکی است. از این دیدگاه، تناقض‌نمای
  باناخ--تارسکی نه‌چندان گزاره‌ای دربارهٔ $\Real^3$، بلکه گزاره‌ای
  دربارهٔ پیچیدگی گروه [انتقال‌ها و دوران‌ها در $\Real^3$] است.
  \cite[p.~16]{Weston2003}
\end{quote}
یعنی اینکه بتوانیم تجزیه‌ای \emph{متناهی} (مانند باناخ--تارسکی) یا
تجزیه‌ای \emph{نامتناهیِ شمارا} (مانند تناقض‌نمای ویتالی) عرضه کنیم،
به واقعیت‌های معینی در نظریهٔ گروه دربارهٔ کار در دو یا سه بُعد
بستگی دارد.

باید پذیرفت که این مشاهدهٔ آخر لطیفهٔ پایان
\olref[banach]{sec} را اندکی خراب می‌کند. چون
«Banach-Tarski» دوبعدی است، اگر بخواهیم قطعاتش را بازآرایی کنیم و
«Banach-Tarski Banach-Tarski» بسازیم، باید آن را به
\emph{بی‌نهایتِ} شمارایی قطعه تقسیم کنیم. برای اصلاح لطیفه باید در
سه بُعد نوشت. این کار را به‌عنوان تمرین به خواننده می‌سپاریم.

یک نکتهٔ پایانی. در \olref[banach]{sec} گفتیم «قطعه‌های» حاصل از
کره نمی‌توانند \emph{اندازه‌پذیر} باشند، بلکه باید
«پراکندگی‌های نامتناهیِ» تجسم‌ناپذیری باشند. دربارهٔ کاربرد انتخاب
در به‌دست‌آوردن \olref{pseudobanachtarski} نیز چنین است. شایسته است
این مطلب را توضیح دهیم.

باز باید اندکی پیش‌زمینه ترسیم کنیم (اما این \emph{فقط} طرحی کلی
است؛ شاید بخواهید مدخلی درسی دربارهٔ \emph{اندازه} را ببینید).
تعریف یک اندازه برای مجموعهٔ $X$ یعنی نسبت‌دادن مقداری
$\mu(E) \in \Real$ به هر $E$ در یک «جبرِ~$\sigma$» روی $X$.
جزئیات در اینجا ضروری نیست، مگر این اصل جمع‌پذیری شمارا: اندازهٔ
اجتماع شمارایی از مجموعه‌های مجزا برابر با مجموع اندازه‌های جداگانهٔ
آن‌هاست؛ یعنی هرگاه $X_n$ها مجزا باشند،
$\mu(\bigcup_{n < \omega} X_n) =
\sum_{n < \omega}\mu(X_n)$. گفتن اینکه مجموعه‌ای «اندازه‌ناپذیر»
است یعنی هیچ اندازه‌ای را نمی‌توان به‌گونه‌ای مناسب به آن نسبت داد.
اکنون با استفاده از $\rotationsgroup$ پیشین:

\begin{cor}[ویتالی]
فرض کنید $\mu$ اندازه‌ای باشد به‌طوری‌که
$\mu(\onesphere) = 1$، و هرگاه $X$ و $Y$ هم‌نهشت باشند،
$\mu(X) = \mu(Y)$. آن‌گاه برای هر
$\rho \in \rotationsgroup$، مجموعهٔ $\funimage{\rho}{C}$
اندازه‌ناپذیر است.
\end{cor}

\begin{proof}
برای برهان خلف، خلافش را فرض کنید. پس برای یک
$\sigma \in \rotationsgroup$ و یک $r \in \Real$ بگذارید
$\mu(\funimage{\sigma}{C}) = r$. بنا بر نامنفی‌بودن اندازه‌ها،
$r \geq 0$. برای هر $\rho \in \rotationsgroup$،
$\funimage{\rho}{C}$ و $\funimage{\sigma}{C}$ هم‌نهشت‌اند؛ پس برای
هر $\rho \in \rotationsgroup$ داریم
$\mu(\funimage{\rho}{C}) = r$. بنا بر \olref{vitalicover} و
\olref{vitalinooverlap}،
$\onesphere =
\bigcup_{\rho \in \rotationsgroup}\funimage{\rho}{C}$ اجتماعی
شمارا از مجموعه‌های دوبه‌دو مجزاست. پس جمع‌پذیری شمارا حکم می‌کند
که $\mu(\onesphere) = 1$ مجموع اندازه‌های همهٔ
$\funimage{\rho}{C}$ها باشد؛ یعنی:
\[
	1 = \mu(\onesphere) = \sum_{\rho \in \rotationsgroup}\mu(\funimage{\rho}{C}) = \sum_{\rho \in \rotationsgroup}r
\]
اما اگر $r = 0$، آن‌گاه
$\sum_{\rho \in \rotationsgroup}r = 0$؛ و اگر $r > 0$، آن‌گاه
$\sum_{\rho \in \rotationsgroup}r = \infty$.
\end{proof}

\end{document}
